\documentclass[12pt]{article}
\usepackage[utf8]{inputenc}
\usepackage[T1]{fontenc}
\usepackage{hyperref}
\title{Tepo Krumen Language Resources for K-12 Curriculum Development}
\author{Compiled by Claude Sonnet 5 (claude-sonnet-5) for Andr\'as Kornai}
\date{}
\begin{document}
\maketitle
\section{Introduction}
Tepo Krumen (Southern Kroumen; ISO 639-3 ted) is a Kru language of south-western C\^ote d'Ivoire with several varieties (Bapo, Dapo, Glawlo, Honpo, Tepo, Wlopo, Yrepo). Search results indicate it is used as a first language but is not taught in schools. Very little educational material was found.
\section{Primary Educational Resources}
None found.
\section{Linguistic and Reference Materials}
A numerals page from the Max Planck EVA numerals database: \url{https://lingweb.eva.mpg.de/channumerals/Krumen-Tepo.htm}.
\section{Religious and Community Resources}
Global Recordings Network lists gospel recordings and notes a New Testament only: \url{https://globalrecordings.net/en/language/ted}, with variety pages \url{https://globalrecordings.net/en/language/12367} (Glawlo) and \url{https://globalrecordings.net/en/language/12365} (Bapo). Joshua Project people-group profile: \url{https://www.joshuaproject.net/index.php/people_groups/11250/IV}.
\section{Digital and Audio Resources}
Audio recordings via Global Recordings Network (above); content is evangelistic.
\section{Curriculum Development Context}
No dictionary, grammar, or school materials were found. Any curriculum would need substantial new documentation work.
\end{document}